\documentclass[11pt,a4paper]{article}

%% PACHETE MATEMATICE
\usepackage{amsmath,amssymb,amsfonts}
\usepackage{amsthm}
\usepackage{mathpazo}
\usepackage{eulervm}
\usepackage{enumerate}
\usepackage{color}
\usepackage{graphicx}
\usepackage[all]{xy}
\usepackage{xcolor}
\usepackage{framed}
\usepackage[unicode]{hyperref}
\setcounter{MaxMatrixCols}{20}
\usepackage{multicol}
%\usepackage[romanian]{babel}


%% ASPECT
\linespread{1.05}
\usepackage[T1]{fontenc}
\usepackage[utf8x]{inputenc}
\usepackage[margin=2cm]{geometry}
% \usepackage[color]{showkeys}
\usepackage{anyfontsize}
\usepackage{titlesec}
\titleformat*{\section}{\large\bfseries}

\usepackage{fancyhdr}
\pagestyle{fancy}
\rhead{\small \color{gray} Notițe de Adrian Manea}
\lhead{\small \color{gray} IntProb-FIA, 2017---2018, L. Lemnete \& M. Olteanu}


\usepackage[autostyle = false, style = german]{csquotes}
\usepackage[romanian]{babel}
\newcommand\qq{\enquote}

%% COMENZILE MELE
\renewcommand*{\proofname}{Dem.:}
\newcommand\bb{\mathbb}
\newcommand\dr{\mathrm}
\newcommand\kal{\mathcal}
\newcommand\fk{\mathfrak}
\newcommand\op{\oplus}
\newcommand\ot{\otimes}
\newcommand\ds{\displaystyle}
\newcommand\id{\indent}
\newcommand\nid{\noindent}
\newcommand\seq{\subseteq}
\newcommand\sneq{\subsetneq}
\newcommand\speq{\supseteq}
\newcommand\spneq{\supsetneq}
\newcommand\rar{\rightarrow}
\newcommand\Rar{\Rightarrow}
\newcommand\xrar{\xrightarrow}
\newcommand\lar{\leftarrow}
\newcommand\Lar{\Leftarrow}
\newcommand\xlar{\xleftarrow}
\newcommand\lrar{\leftrightarrow}
\newcommand\Lrar{\Leftrightarrow}
\newcommand\ideal{\trianglelefteq}
\newcommand\ai{\text{ a.\^{i}. }}
\newcommand\sk{(\textit{Schi\c{t}\u{a}}) }
\newcommand\rhup{\rightharpoonup}
\newcommand\rhdn{\rightharpoondown}
\newcommand\lhup{\leftharpoonup}
\newcommand\lhdn{\leftharpoondown}
\newcommand\vid{\emptyset}
\newcommand\wh{\widehat}
\newcommand\ol{\overline}
\newcommand\NN{\mathbb{N}}
\newcommand\RR{\mathbb{R}}
\newcommand\ZZ{\mathbb{Z}}
\newcommand\QQ{\mathbb{Q}}
\newcommand\CC{\mathbb{C}}

%% STILURILE MELE
\newtheoremstyle{dotless}{}{}{\itshape}{}{\bfseries}{:}{ }{}

\theoremstyle{dotless}
\newtheorem{theorem}{Teorem\u{a}}[section]
\newtheorem{proposition}{Propozi\c{t}ie}[section]
\newtheorem{lemma}{Lem\u{a}}[section]
\newtheorem{corollary}{Corolar}[section]
\newtheorem{exercise}{Exerci\c{t}iu}

\newtheoremstyle{dotlessdr}{}{}{}{}{\bfseries}{:}{ }{}
\theoremstyle{dotlessdr}
\newtheorem{example}{Exemplu}[section]
\newtheorem{definition}{Defini\c{t}ie}[section]
\newtheorem{remark}{Observa\c{t}ie}[section]




\begin{document}

\begin{center}
  {\large\textbf{Seminar 2 \\ Integrale improprii cu parametri}}
\end{center}

\vspace{1cm}

1. Să se calculeze integralele, folosind derivarea sub integrală:
\begin{enumerate}[(a)]
\item $ I(m) = \ds\int_0^{\frac{\pi}{2}} \ln(\cos^2 x + m^2 \sin^2 x) dx, m > 0 $;
\item $ I(a) = \ds\int_0^{\frac{\pi}{2}} \ln(a^2 - \sin^2 \theta) d\theta, a > 1 $;
\item $ I(a) = \ds\int_0^1 \frac{\arctan ax}{x \sqrt{1 - x^2}} dx, a > 0 $.
\end{enumerate}

\emph{Soluții:}

(a) Dacă considerăm funcția:
\[
  f(x, m) = \ln(\cos^2 x + m^2 \sin^2 x),
\]
observăm că este continuă și admite o derivată parțială continuă în raport cu $m$.

Atunci obținem:
\[
  \frac{\partial f}{\partial m} = I'(m) = 2m \int_0^{\frac{\pi}{2}} %
  \frac{\sin^2 x}{\cos^2 x + m^2 \sin^2 x} dx
\]

Pentru a calcula integrala, facem schimbarea de variabilă $\tan x = t$ și atunci:
\[
  dt = \frac{1}{\cos^2 x} dx.
\]

Integrala inițială se poate prelucra:
\[
  I'(m) = 2m \int_0^{\frac{\pi}{2}} \frac{\sin^2 x}{cos^2 x(1 + m^2 \tan^2 x)} dx.
\]
Așadar, pentru a obține $\sin^2 x$ în funcție de $t$ calculăm:
\[
  \frac{\sin^2 x}{1 - \sin^2 x} = t^2 \Rightarrow \sin^2 x = \frac{t^2}{1 + t^2}.
\]

În fine, integrala devine:
\[
  I'(m) = 2m \int_0^\infty \frac{t^2}{(1 + m^2 t^2)(1 + t^2)} dt.
\]
Făcînd descompunerea în fracții simple, obținem:
\[
  \frac{t^2}{(1 + m^2 t^2)(1 + t^2)} = \frac{1}{m^2 - 1} \Big( \frac{1}{t^2 + 1} - %
  \frac{1}{m^2 t^2 + 1} \Big).
\]
și calculăm în fine integrala $I'(m) = \frac{\pi}{m + 1}$. Integrăm și găsim
$I(m) = \pi \ln(m + 1) + c$. Deoarece $I(1) = 0$, rezultă $c = -\pi \ln 2$
și, în fine:
\[
  I(m) = \pi \ln \frac{m + 1}{2}.
\]

\vspace{1cm}

(b) Dacă considerăm funcția:
\[
  f(\theta, a) = \ln(a^2 - \sin^2 \theta),
\]
observăm că este continuă și admite o derivată parțială continuă în raport cu $a$.

Atunci avem:
\[
  I'(a) = \frac{\partial f}{\partial a} = \ds\int_0^{\frac{\pi}{2}} \frac{2a}{a^2 - \sin^2 \theta} d\theta.
\]

Cu schimbarea de variabilă $t = \tan \theta $, avem succesiv:
\begin{align*}
  I'(a) &= \int_0^\infty \frac{2a}{a^2 - \frac{t^2}{1 + t^2}} \cdot \frac{1}{1 + t^2} dt \\
        &= \int_0^\infty \frac{2a}{t^2(a^2 - 1) + a^2} dt \\
        &= \frac{2a}{a^2 - 1} \int_0^\infty \frac{1}{t^2 + \frac{a^2}{a^2 - 1}} dt \\
        &= \frac{2a}{a^2 - 1} \cdot \frac{\sqrt{a^2 - 1}}{a} \cdot \arctan t \frac{\sqrt{a^2 - 1}}{a} \Big|_0^\infty \\
        &= \frac{2}{\sqrt{a^2 - 1}} \cdot \frac{\pi}{2} \\
        &= \frac{\pi}{\sqrt{a^2 - 1}}.
\end{align*}

Integrăm pentru a obține $I(a)$ și găsim:
\[
  I(a) = \int \frac{\pi}{\sqrt{a^2 - 1}} da = \pi \ln(a + \sqrt{a^2 - 1}) + c.
\]

Pentru a calcula constanta $c$, putem rescrie integrala din forma inițială:
\begin{align*}
  I(a) &= \int_0^{\frac{\pi}{2}} \ln\Big( a^2\Big(1 - \frac{\sin^2 \theta}{a^2}\Big) \Big) d\theta \\
       &= \int_0^{\frac{\pi}{2}} \ln a^2 d\theta + \int_0^{\frac{\pi}{2}} \ln \Big(1 - \frac{\sin^2 \theta}{a^2} \Big) d\theta.
\end{align*}

Putem considera limita $ a \to \infty $ și atunci integrala de calculat tinde la 0, deci $ c = -\pi \ln 2 $.

Concluzie: $ I(a) = \pi \ln \frac{a + \sqrt{a^2 - 1}}{2} $.

\vspace{1cm}

(c) Dacă considerăm funcția $ f(x, a) = \frac{\arctan ax}{x \sqrt{1 - x^2}}$, observăm că este continuă
și admite o derivată parțială continuă în raport cu $ a $. Atunci:
\[
  I'(a) = \frac{\partial f}{\partial a} = \int_0^1 \frac{dx}{(1 + a^2 x^2) \sqrt{1 - x^2}}.
\]

Pentru a calcula integrala, facem schimbarea de variabilă $ x = \sin t $ și obținem:
\[
  I'(a) = \int_0^{\frac{\pi}{2}} \frac{dt}{1 + a^2 \sin^2 t}.
\]

Mai departe, aplicăm schimbarea de variabilă $ \tan t = u $ și obținem succesiv:
\begin{align*}
  I'(a) &= \int_0^\infty \frac{du}{(1 + a^2)u^2 + 1} \\
        &= \frac{\pi}{2} \cdot \frac{1}{\sqrt{1 + a^2}}.
\end{align*}

Așadar, printr-o integrare în funcție de $ a $, găsim:
\[
  I(a) = \frac{\pi}{2} \ln(a + \sqrt{1 + a^2}) + c.
\]

Cum $ I(0) = 0 $, găsim $ c = 0 $.

\vspace{2cm}

2. Să se calculeze, folosind funcțiile $ B $ și $ \Gamma $:
\begin{enumerate}[(a)]
\item $ \ds\int_1^e \frac{1}{x} \ln^3 x \cdot (1 - \ln x)^4 dx $;
\item $ \ds\int_0^1 \sqrt{\ln \frac{1}{x}} dx $;
\item $ \ds\int_0^\infty \frac{x^2}{(1 + x^4)^2} dx $;
\item $ \ds\int_{-1}^\infty e^{-x^2 - 2x + 3} dx $;
\end{enumerate}

Indicație (d):
\[
  \int_{-1}^\infty e^{-(x + 1)^2 + 4} dx = e^4 \int_{-1}^\infty e^{-(x + 1)^2 } dx.
\]
Notînd acum $ x + 1 = y $, integrala se rezolvă cu formula pentru $ \Gamma $.

\end{document}